\section{Related Work}\label{sec:related}

Most testing frameworks work for Solidity smart contracts
for Ethereum-like blockchains, that still cover a majority of the
running smart contracts.
Pytest~\citep{pytest} is a Python test suite for Ethereum smart contracts.
Foundry~\citep{foundry} is a Rust framework
that includes a testing suite
where Solidity contracts create the fixture and runs the tests.
Remix includes a similar test suite in Solidity~\citep{remixunittesting},
integrated inside the development environment.
Both Waffle~\citep{waffle} and Hardhat~\citep{hardhat} are TypeScript
testing frameworks for Solidity smart contracts:
they use the Web3 API tor install the contracts in blockchain
and interact with them.
The legacy Truffle framework~\citep{truffle} is now discontinued
and converging into Hardhat.
They are probably the most similar approach to our use of JUnit,
despite the different programming language.
The Ape framework is a Web3 development tool for different smart contract
languages, including Solidity. It has its own testing
framework~\citep{apetesting}, based on pytest.
Wake~\citep{wake} is a development and testing framework based on pytest,
that includes some fuzzing feature.

Other blockchains, with different smart contract languages, have their own
testing frameworks. We cannot cite them all here.
Just for providing a few examples, there are testing frameworks
for Solana~\citep{rapidinnovation}, Polkadot~\citep{zombienet},
and Algorand~\citep{algokit}.

Writing testcases for smart contracts is considered difficult
since synamic sequences of behaviors are tested, rather than
specific outputs. Therefore, the \emph{test oracle problem} plays
a significant role here and it is why metamorphic testing
(that tests consistency in the output rather than the exact output)
is widely used for testing smart contracts~\citep{SunJZL25,VillanuevaSR25}.
Moreover, testcases for smart contracts can
sometime be generated automatically,
for instance for specific properties such as gas consumption
bounds~\citep{AlbertDDFG26} or with the support of large language models
and retrieval augmented generation~\citep{SbouiL26}.
Fuzzing has also been used to generate or strengthen testcases for
smart contracts~\cite{LahamiKMMR24,WongZNLCYLLW24}.

Testcases can be tested themselves, in order to evaluate their
efficacy and adequacy. This is typically achieved through
mutation-based techniques~\citep{ChenZZZW24,BanescuB0PZ24,Barboni0PBZ24}.

Finally, a systematic review on smart contract testing is reported
in~\citep{ImperiusA22}.

This article builds on the JUnit testing framework~\citep{junit},
which is the de facto standard for Java. Namely, we test smart contracts
written in a subset of Java called Takamaka~\citep{Spoto19} and
executed over the public blockchain Hotmoka~\citep{hotmoka}. The latter
is normally based on the proof of space consensus engine Mokamint,
but can also run over the Tendermint engine~\citep{Kwon14},
leading to a proof of stake version of Hotmoka.
